\documentclass[12pt]{article}
\usepackage[utf8]{inputenc}
\usepackage[T1]{fontenc}
\usepackage{hyperref}
\title{Paloor Language Resources for K-12 Curriculum Development}
\author{Compiled by Claude Sonnet 5.5 (claude-sonnet-5-5) for Andr\'as Kornai}
\date{}
\begin{document}
\maketitle

\section{Introduction}
Paloor (also Palor, Falor; ISO 639-3 \texttt{fap}, Glottocode palo1243) is a Cangin language of Senegal, spoken by the Palor people west of Thi\`es. Wikipedia gives about 10,700 speakers (2007); ASJP lists 22,000. Wikitongues lists it as EGIDS 6b (threatened). Its closest relative is Ndut. The searches found only descriptive pages; no educational materials.

\section{Primary Educational Resources}
None found.

\section{Linguistic and Reference Materials}
\begin{itemize}
\item Wikipedia, ``Palor language''. \url{https://en.wikipedia.org/wiki/Palor_language}
\item Wikipedia, ``Cangin languages''. \url{https://en.wikipedia.org/wiki/Cangin_languages}
\item ASJP wordlist for Palor. \url{https://asjp.clld.org/languages/PALOR_3}
\end{itemize}

\section{Religious and Community Resources}
\begin{itemize}
\item Joshua Project profile, ``Jola, Falor'' in Senegal. \url{https://joshuaproject.net/people_groups/11185/SG}
\item Wikipedia, ``Palor people''. \url{https://en.wikipedia.org/wiki/Palor_people}
\end{itemize}

\section{Digital and Audio Resources}
\begin{itemize}
\item Wikitongues language page for Paloor (fap). \url{https://wikitongues.org/languages/fap/}
\end{itemize}

\section{Curriculum Development Context}
No teaching materials, orthography or scripture were found. Curriculum work would need to start from documentation and community involvement; the closely related Ndut and Cangin literature may offer comparative help.

\end{document}